\documentclass[12pt]{article}

\usepackage[a4paper,margin=1in]{geometry}
\usepackage{amsmath, amssymb}
\usepackage{setspace}
\usepackage{titlesec}

\setstretch{1.15}

\title{Lecture Scribe — CSE400: Fundamentals of Probability in Computing\\
Lecture 3: Introduction to Probability Theory}
\author{Instructor: Dr. Dhaval Patel}
\date{Date: January 13, 2026}

\begin{document}

\maketitle

\section{Course Context and Lecture Scope}

\subsection{Course Identification}

Course Title: CSE400 – Fundamentals of Probability in Computing

Lecture Topic: Introduction to Probability Theory

Instructor: Dr. Dhaval Patel

Institution: SEAS, Ahmedabad University

This lecture introduces:

General course information

Motivation for studying the course

Engineering applications

Course structure and participation framework

(As presented in slides 1 and Outline slides) 


\section{Definitions and Notation (From Lecture Slides)}

The lecture slides for this session do not introduce formal mathematical definitions, notation, probability axioms, or symbolic representations.

Therefore:

No definitions of probability, random variables, sample spaces, or events are stated in this lecture.

No formal notation is introduced in this lecture.

This lecture functions as an introductory and organizational lecture rather than a theoretical one.

\section{Assumptions and Conditions}

The following assumptions and conditions are explicitly stated or implied through course structure slides:

\subsection{Learning Assumptions}

Active participation is expected through Campuswire.

Learning includes discussion, feedback, and collaborative engagement.

Students are expected to participate in active learning and class discussions.

(From “Active Learning and Class Discussion” slides) 


\subsection{Communication Conditions}

Queries should preferably be posted on Campuswire.

Direct messaging with instructor is allowed for private discussions.

Contact hours are available through Campuswire.

(From “Timing for Discussion and Difficulty Sessions” slide) 


\section{Statements of Results or Theoretical Content}

This lecture does not contain:

Theorems

Lemmas

Propositions

Mathematical results

Proof statements

Hence, no formal statements or results are included in this lecture.

\section{Step-by-Step Proofs}

No proofs are presented in this lecture.

\section{Worked Examples (From Lecture Slides)}

The lecture includes contextual examples rather than mathematical worked examples.

Example 1: Motivation for Learning CSE400

From the outline slide:

Reason for learning the course:

Example: Daily life conversations

Step representation (as shown in slides):

Probability concepts arise in daily life situations.

Conversations involving uncertainty motivate learning probability.

(No mathematical steps are presented in the slides.) 


Example 2: Engineering Applications

Engineering applications listed in slides:

Speech Recognition

System Radar System

Communication Network

Step representation (as shown in slides):

Probability concepts are applied in engineering domains.

These applications motivate study of probability in computing systems.

(No computational or numerical example is shown.) 


\section{Course Structure Relevant for Examination Preparation}

\subsection{Active Learning Platform}

Campuswire is used for:

Anonymous participation

Posting and back-channel communication

Real-time feedback and polling

Direct communication with instructor/TAs

(From Campuswire slide) 


\subsection{Lecture Schedule}

Lecture sessions:

Section 1:

9:30 AM – 11:00 AM (Tuesday, Thursday)

GICT Room 136

Section 2:

1:00 PM – 2:30 PM (Tuesday, Thursday)

GICT Room 137

(From schedule slide) 


\subsection{Project Component (Overview)}

Project weightage:

30\% of course

Major milestones:

Concept Evolution Maps

Mathematical Modeling

Coding and Simulation

Inference and Randomized Algorithms

Application of Randomized Algorithms

Bounds, Analysis, and Final Deliverables

(From project execution slides) 


\section{Summary for Revision}

From this lecture:

The lecture introduces course logistics and motivation.

No mathematical probability theory content is formally introduced.

No definitions, notation, theorems, or proofs are included.

Motivation is provided through:

Daily life conversations

Engineering applications such as speech recognition, radar systems, and communication networks.

Active learning and structured participation form part of the course framework.

\end{document}
